\section{서론}

한 경주의 단승식, 복승식, 쌍승식, 삼복승식과 삼쌍승식은 같은 실현결과에
연동되지만 서로 다른 사건에 지급하고 별도 패리뮤추얼 풀에서 청산된다. 따라서
각 풀의 마감 배당률은 동일한 기초 순위에 대한 서로 다른 가격측도를 제공한다.
이 구조는 두 질문을 낳는다. 여러 풀의 전체 가격벡터는 하나의 상위 3위
상태공간에서 나온 주변분포와 얼마나 가까운가? 그 정합성에서 벗어난 가격차를
복합복권의 행동모형이 강화된 순위 기준모형보다 더 잘 설명하는가?

기존 경마시장 연구는 주로 단승 배당률이 시사하는 확률과 사후 승률을 비교하여
인기마--비인기마 편향을 검정했다. 비인기마가 과대 베팅되고 인기마가 과소
베팅된다는 결과는 오래된 경험적 규칙성이지만
\citep{griffith1949,weitzman1965,ali1977}, 홍콩처럼 편향이 약하거나 반대로
나타난 시장도 있다 \citep{buschehall1988}. 한국 자료에서는 말과 경주의
이질성을 허용한 승률모형으로 전형적 편향이 보고됐고
\citep{jeongetal2019}, 게시 배당률의 반올림 자체가 전통적 검정통계를
왜곡할 수 있다는 지적도 있다 \citep{wallsbusche2003}. 이 문헌은 위험선호,
확률 오인, 참가자 이질성과 측정오차를 구분해 왔지만, 한 시장의 배당률과
실현결과만으로는 위험선호와 확률오인이 같은 가격패턴을 만들 수 있다.

\citet{snowbergwolfers2010}는 단승식에서 추정한 두 설명이 쌍승·복승·삼쌍승
가격으로 이동하는지를 비교하여 이 관측적 동등성에 제약을 가했다. 다만 이들의
자료에는 exotic 승식의 당첨 배당만 있고 비당첨 조합의 가격은 없었다.
\citet{koivurantakorhonen2019}은 핀란드와 스웨덴의 완전한 exotic 가격을
이용하여 당첨조합 선택 문제를 줄이고, 1위와 2위를 순차적으로 평가하는 모형을
검정했다. 따라서 완전한 복합승식 가격 자체는 새로운 기여가 아니다. 남은
질문은 순서 있는 상위 3위 전체 상태를 여러 하위 승식과 동시에 연결할 때,
교차풀 가격 정합성과 세 번째 조건부 평가단계를 무엇까지 식별할 수 있는가이다.

이 논문은 연구를 두 층으로 나눈다. 첫째 층인 주분석에서는 삼쌍승식의 각
조합을 ``1위 $i$, 2위 $j$, 3위 $k$''라는 원자적 상태로 해석한다. 삼쌍승식
역배당률을 경주 안에서 정규화한 뒤 합산하면 단승, 쌍승, 복승, 삼복승이
지급하는 사건의 재구성 가격을 얻는다. 이 가격을 각 승식의 실제 정규화 가격과
비교하여 객관적 승률을 추정하지 않고도 교차풀 정합성을 측정한다.

게시 표시상한은 이 비교의 부차적 문제가 아니다. 상한 없는 경주의 점추정은
정확하지만 표본선택에 노출되고, 상한 경주를 포함한 전체표본은 가격을 구간으로만
식별한다. 따라서 이 논문은 clean 표본의 점추정과 전체표본의 부분식별 경계를
공동 주결과로 두고, 두 결과가 같은 방향을 지지할 때만 전체 시장에 관한 강한
결론을 허용한다.

둘째 층은 Snowberg and Wolfers의 검정논리를 세 번째 순위까지 확장하되,
실현 착순으로 조건부 순위확률을 표본 밖에서 추정한 뒤, 단승 가격에 맞춘
확률가중 함수가 다른 풀의 전체 가격벡터를 설명하는지 검정한다.
특히 삼쌍승 사건의 결합확률에 확률가중을 한 번 적용하는 축약 모형,
`1위→나머지 두 순위`의 2단계 모형, `1위→2위→3위`의 3단계 비축약 모형을
비교한다. 순서가 없는 승식에는 임의의 사건트리를 부여하지 않고, 순서형 모형의
풀 밖 예측을 검증하는 데 사용한다. 동시에 게시가격 Harville(비보정)과 연도 밖
착순으로 추정한 보정·단계조정 Harville을 같은 경주에 적용하여 행동모형의 추가 설명력을
직접 비교한다.

결과는 두 층의 구분이 실제로 중요함을 보여준다. 주분석에서 삼쌍승 주변가격은
하위 복합승식을 Harville과 비정보적 기준보다 일관되게 잘 재구성하지만,
사전에 정한 강한 절대 TV 기준은 clean 점표본과 전체표본 경계를 동시에
통과하지 못한다. 진단분석에서 보정·단계조정 조건부 순위확률은 보정 Harville($\alpha=1$)보다
2·3위의 연도 밖 로그손실과 보정오차를 낮춘다. M-S2/M-S3은 M-R보다
가격오차가 작지만, 동일표본의 단계조정 Harville을 네 승식 어디에서도 넘지
못한다. 이 판정은 공통지지 조합과 최저가격 10\%를 제외한 조합에서도 유지된다.
따라서 가격의 교차풀 정합성은 높지만 완전하지 않으며, 그 잔여 차이를 비축약
행동으로 구조적으로 해석할 증거는 얻지 못한다.

이 연구의 기여는 세 가지다. 첫째, 완전한 삼쌍승 가격벡터와 고정된 주변화
행렬을 이용하여 다섯 풀의 가격을 하나의 상위 3위 상태공간에서 비교한다.
둘째, Harville·동일 경주 순열·타경주·균등 기준을 같은 표본과 손실함수에서
비교하고, 점추정과 전체표본 부분식별 경계를 나란히 두어 높은 적합도가 공통
인기순위·차원 또는 선택된 clean 표본의 산물인지 구분한다. 셋째,
복합복권의 축약 여부를 세 번째 순위까지 확장하되 게시가격 Harville(비보정)과
보정·단계조정 Harville을
동일표본 기준으로 함께 두어, 내부적인 M-R 대비 개선과 독립적인 행동설명력을
구분한다.

여기서 ``재구성력이 높다''는 말은 삼쌍승식 가격이 객관적 확률이라는 뜻이
아니다. 별도 풀의 정규화 역배당률을 같은 가격측도의 주변분포로 해석하려면
정보시점, 공제·반올림, 참가자 구성과 평가기능에 관한 유지가정이 필요하다.
그 가정을 유지하지 않으면 주분석은 서로 다른 풀 가격측도의 정적 정렬 정도를
측정한다. 또한 M-S의 M-R 대비 우위는 집계가격의 내부 모형순위일 뿐이며,
더 강한 비행동 기준모형을 이기지 못한 이상 개별 베팅자의 심리과정에 대한
증거가 아니다. 본문은 이 두 해석의 경계를 분리한다.
